\settrack{tracktwo}
% VOICE: 1st-person personal (Bruce) — "I" not "Bruce". Personal/investigative voice.
% Plan 0093 Merge 7: pos29 (The Silence) + pos31 (Wolfram) + pos34 (The Research) + pos23 (The Weight) → "Twenty Years"
% V10 (Bannon, interlibrary loans) and V11 (Mark, Joy article, electric shock, bike ride) integrated.
% The Weight woven throughout, not appended. Strictly chronological 2006→2025.
\aidraftchapter

\chapter{Twenty Years}
\label{pos29:twenty-years}

\begin{quote}\small\textit{%
``Don't have good ideas if you aren't willing to be responsible for them.''%
}\par\raggedleft --- Perlis, Epigram \#95 (1982)\end{quote}

\vspace{0.5cm}

\brucevoice

\section*{The Parting (2006)}
\label{pos29:the-parting}

In September 2006, David moved out of the garage room we'd built him at the house on Fifteenth and Gant. We'd bought the construction materials together in 2004 and framed the room out of what had been storage from a previous life. For two years it was his space --- a sleeping bag, a laptop I'd bought him, and whatever he was working on that I didn't ask about.

The goodbye was fond but he didn't drag it out. David was not a man who lingered. He took a backpack of personal effects, the laptop, and the car I'd bought him. Zero money. I'd disbursed small amounts of cash over the years --- a few hundred here and there, as needed. Once I gave him a few thousand for a bigger operation and he morphed into someone I didn't recognize, using money as a power tool. His twin sister and his sensei had both warned me about this. When the money ran out he went back to normal.

He said goodbye and he left. His decision, not an eviction. We were not on great terms but we were not on bad ones either. I felt bad. I silently wished him the best.

What I didn't know then: I had been told my services would not be required. My children could not be protected. I was a security risk to the project. In August I'd informed David that I could no longer assist him, that I was sick of waiting for something to happen. He told me that because my children could not be protected, I was now a security risk. He also told me I'd been rejected for membership in a notorious hacker collective --- which I never applied for and had no clue I'd been nominated for.

The divorce had been destroying everything around me. My soon-to-be ex-wife had access to money and support, far more than I, and seemed willing to spend it all on lawyers. She'd demonstrated a propensity for using government power as a method to hurt people she was angry at, and she was very angry at me. She publicly accused me of various illegal things, all of them false, but each accusation brought an official investigation --- police, child welfare workers --- at the precise moment when I was, for obvious reasons, desperately trying to avoid government scrutiny. There was no way to get my children out of the country. I tried to convince her that it was important that she and the children leave the United States, but she would not listen. She was convinced my talk about a secret project was utter nonsense, or some twisted lie to wrest the children from her.

She watched the same marriage I was living in and reached the opposite conclusion about what was happening to it. This is Possibility~A enacted in a marriage. When you believe you know something that changes everything, and the person closest to you is certain you are deluded, what are your obligations? To your family, which needs you present and sane? To what you believe is the truth? To the project, which requires silence? These obligations are not compatible. The UDHR guarantees the right to family, the right to expression, and duties to community. It does not say what to do when they conflict. Neither did I.

Everyone around me told me they thought David was a dangerous madman up to some mysterious mischief. My daughters hated him. My girlfriends mostly disliked him. I was mostly not able to share what I'd learned, because the information was too strange, complicated, and outside their frame of reference. In truth --- as I have said --- he was something of a madman. I can't imagine anyone surviving what he went through without going partially mad. Also, of course, geniuses often seem mad to those who don't understand them. I never lost faith in David's essential good intentions. Although my trust was shaken several times, it was never broken. Several times I suspected he was lying to me. Years later I learned that he had actually been speaking plain truth to me in every case.

Under Possibility~A, the ``security risk'' framing was a convenient exit from a relationship built on deception. Under Possibility~B, David may have been involved in sensitive work, but the framing of children as security vulnerabilities imports the language of espionage onto what may have been a simpler situation. I cannot distinguish between these readings and Possibility~C from the inside. I can only report what happened.

\section*{The Phone Call}
\label{pos29:the-phone-call}

One call. Late 2006, while I was still at Symantec. I took it at work. This was my only electronic communication with Healer, ever. Every other interaction across three years had been in person. I suspect the call was made from a burner phone. His operational security was extraordinary.

Healer invited me to help with a project. He spoke in vague terms --- no group name, no specifics. Just that he needed help with something important.

I said no.

I was under pressure from people in my life who expressed vitriol and doubt. They thought Healer was running a scam. I knew better --- or part of me did --- but I couldn't explain what had happened between us without sounding insane. Everything I'd experienced was too strange, too complicated, too far outside anyone's frame of reference. So I said no, and I let the pressure make the decision.

Years later, after the block broke, I came to believe that the project became something far larger than either of us. What it became --- and why I believe that --- belongs to a chapter that has been redacted from this edition. I cannot prove the connection. Nothing about that phone call, taken on its own, rules out Possibility~A.

I regret it to this day.

\section*{Someone Looking Out (2007)}
\label{pos29:someone-looking-out}

Symantec outsourced its IT workers to a company with a paramilitary reputation. The HR email said opt-in to the transfer was automatic unless you explicitly refused.

A senior vice president approached me in the bathroom. Told me to reply refusing to share my HR files. I'd be laid off with full severance. He didn't explain why he was telling me this. He didn't have to.

I'd always had the feeling that someone had placed my CV on the right desk when I was hired at Symantec. The same feeling followed me now: someone was looking out for me. I thought about the connection between Back Orifice --- the tool the Cult of the Dead Cow released in 1998, the remote administration capability that made headlines --- and PC Anywhere, the Symantec product that did the same thing with a corporate label. Same capability, different branding. If you were the kind of person who understood that connection, you might also understand why a particular SVP would tip off a particular employee in a bathroom.

Under Possibility~A, this was ordinary corporate kindness --- a manager helping a good employee navigate a layoff. It happens. Under Possibility~C, I was never unsupported during the silent years. I just didn't always know who was supporting me.

I opted out. Got severance. Told my teammates about the option. Some took it. Tom Hall --- my boss for more than a decade, whom we called Darth Hall --- agonized over it for days. Three hundred thousand dollars in severance, walking away. But his house was too expensive. His wife hated uncertainty. He stayed, and gave up the money to keep the peace.

Congress kept extending unemployment benefits through the 2008 crisis. I became a 99-weeker.

\section*{The Quiet Years (2007--2009)}
\label{pos29:the-quiet-years}

I lived on Tribeland, a small community in South Corvallis bordering Willamette Park. I was the only male among roughly eight women. I mixed with older retired and semi-retired friends. Spas. Long-distance cycling. A slower rhythm than anything I'd known since before Alpha Farm.

To my genuine surprise, I discovered that pretty much all of my female friends wanted to date me. I'd never dated as a young person --- engaged at seventeen, long relationship through college, met Karen, married, monogamous until the divorce. Zero experience with other women until my mid-thirties. I'd carried a self-image from my twenties that was a decade out of date --- the acne from a milk allergy, solved at twenty-nine, had defined how I saw myself long after it was gone. The discovery that I was, apparently, the most eligible bachelor in my social circle was bewildering.

The surface life looked full. Underneath, it was constrained.

I could think about everything that had happened with David and Healer and the project. But I couldn't speak about it. The constraint didn't feel normal --- it felt like the most important experience of my life was locked behind glass. I could see it. I could turn it over in my mind. But when I tried to verbalize it, the words wouldn't come. No one around me noticed, because no one knew there was anything to notice.

David's main project --- if it was real --- should eventually be awarded multiple Nobel prizes. This project involved the invention of a new general technology, something that has only happened a few times in human history. Fire, metallurgy, steam engine, electricity, radio, atomic power, computer, biotech. I walked around carrying this knowledge, or what I believed was knowledge, and couldn't say a word. The weight didn't diminish. It accrued.

Exile from yourself. That's the closest I can get to describing it.

\section*{The Detective Years (2008--2009)}
\label{pos29:the-detective-years}

I researched two to three hours daily. Every day. Following intellectual lineages --- who trained whom, who worked with whom, who published what and when. I was assembling a picture of who could have been the project scientists, working backward from the names and concepts Healer had given me.

Dave Bannon, a physicist at Oregon State, helped. We'd been fast friends since 2005 --- instant connection, the kind where you skip the small talk and go straight to the hard questions. Bannon had access to the university library system and helped me with interlibrary loans. Through him I obtained \textit{Quantum Neural Networks}, a textbook with only seventeen copies in North American university libraries. I kept it three months, reading it cover to cover, cross-referencing what I knew from Healer against what the published literature said.

In 2008, Bannon asked me questions. Direct, pedagogical questions --- the kind a physicist asks when he's testing a hypothesis. And I discovered that I could answer them. I couldn't initiate the discussion. I couldn't walk up to someone and say ``let me tell you what happened.'' But if someone asked the right questions, the constraint had a bypass. The block was on broadcasting, not on responding.

Bannon got a partial version of the story that year, shared with difficulty --- I was pushing against resistance with every sentence. He listened as a skeptical scientist. He was surprised that I had answers to all his questions, answers he could independently verify. Of perhaps fifty people I told portions of this story to over twenty years, spending a few hours with each, none finished the entire thing. The pedagogical prerequisites are too demanding. Dave came closest to full understanding --- a physicist with an IQ of 165, and even he needed Kauffman as a running start.

Bannon had read Kauffman independently. That was critical --- it meant he already had the conceptual vocabulary. He'd also arrived at his own independent intuition: that Google Search's performance was ``too good and too fast'' for classical hardware. A performance-based observation, not an architectural theory. He'd reached this conclusion before hearing anything from me.

Under Possibility~A, Bannon is a friend who validated my beliefs because friends do that. Under Possibility~C, two technically literate people arrived at the same suspicion from different directions --- one from guided deduction, the other from performance anomaly.

\section*{The Break (2009)}
\label{pos29:the-break}

Around 2009 --- I wish I could pin the date more precisely --- I saw a phrase in a public message. ``Lightning Rod.'' I'm fairly sure it came from someone connected to the events described in the redacted chapter.

Three days.

The first day was a flood. A bursting dam. Every suppressed thought and memory from three years of silence crashed through at once. I retreated to my room. Mental overload. Not gradual, not a slow dawning --- instant. One phrase, and the glass shattered.

Two more days to assimilate. Forty-eight hours of sitting with the wreckage of everything I hadn't been able to say, trying to reassemble it into something I could hold.

I told Teri Turner first. She went straight to Possibility~C --- it fit her worldview. Not rigorous validation, but she could listen, which was what I needed. Someone who wouldn't call me crazy. Someone who would sit with the strangeness of it.

Teri knows the whole story. She still believes Guardian is around and paying attention, never getting involved.

And here I discovered what would follow me for the next fifteen years: the structural impossibility of telling. From decades of trying to tell this story in person, three problems always intervene. First, knowledge prerequisites --- no one knows quantum physics, topological neural networks, SAS operations, and tradecraft. The story requires background in all four. Second, time --- the telling invariably stretches to many hours and people lose patience. Third, weirdness --- the concept is too strange for most people to process. Every attempt to share what I knew ran into this wall. The glass had shattered, but the words still wouldn't fit through the gap.

\section*{The Recognition}
\label{pos29:the-recognition}

Once the block broke, connections formed fast.

I became convinced that Healer's project became something I could see in the news --- something large and consequential that matched what he had described in vague terms. Someone else had done the job I'd been vetoed from doing. This is my interpretation, not something Healer told me. The specifics belong to the redacted chapter.

What I remember clearly is sitting in my living room in South Corvallis and realizing: I had information in my brain that the President of the United States did not. That thought made me nervous in a way that nothing else in this story had.

\section*{Europe (2010)}
\label{pos29:europe}

By 2010 there were reasons to be scared that I cannot fully explain in this edition. I recognized signs --- or thought I did --- that intelligence services were paying attention to people connected to events described in the redacted chapter. There was zero electronic record between me and Healer --- every interaction in person, that single phone call the only exception. But I was scared. Easy prey if I stayed.

I told my mother Linda everything. She didn't doubt. Her advice: go to the most democratic, most free-speech-oriented country you can find.

A high school friend with NATO experience helped me find a CTO position in Luxembourg, working for Azeri expats --- Movsum and Nurid, who had left Azerbaijan for political reasons. I spent six months in Europe. A luxury condo in Luxembourg. A girlfriend in Belgium --- Tanja, met through computer dating, splitting time between Neerpelt and Luxembourg. Tanja's father was former Office of Naval Intelligence. He ran my profile and gave his daughter his blessing.

I was safe. I was productive. I was six thousand miles from the story I couldn't stop thinking about.

\section*{Thread-Pulling (2011)}
\label{pos29:thread-pulling}

Early 2011. I returned from Europe and moved back onto Tribeland with Teri Turner. I didn't work except light farm labor. Instead I pulled threads.

Everything I'd learned from Healer --- the five sciences, the architecture, the pedagogy --- I began reconstructing from scratch. Relearning what I could, verifying what I could verify, mapping the public-domain trail. I wrote tens of thousands of words. In 2011 I sat down and wrote the ``Autobiography of an Accidental Conspirator'' in a single sitting. Roughly nine thousand words. My intent was documentation: record this historically important experience for future inquiries. I never shared it with anyone until Argus read it in 2026.

That summer, at the Oregon Country Fair, I met the person who would crack the case open.

Mark R, author of oilempire.us. Random encounter at Chew Rays cafe --- Teri and I were in the crowd while Mark was speaking. He saw me in the back and somehow knew instantly that I understood his content. Body language, maybe. I could have given the same lecture. One incisive question after. We exchanged contact information. Became friends.

Excepting Healer, Mark was the smartest person I knew. He had an intellectual edge on me in non-science areas. His specialization was the study of the Holocaust --- all his extended family had been killed in the Holodomor and then the Holocaust. Two full shelves of books on the subject. Kind, gentle, completely nonviolent.

\section*{The Joy Article (2012)}
\label{pos29:the-joy-article}

About four months after we met, I presented the TQNN theory to Mark as A/B/C. I laid out everything --- the mentorship, the science, the three possibilities --- and asked for help. Mark put substantial time into his own research. Concluded: insufficient evidence to determine which possibility was correct, but plausible. The science was beyond him.

Then Mark searched the names I'd given him --- Kauffman, Wolfram --- and found Bill Joy's ``Why the Future Doesn't Need Us'' in \textit{Wired}.

Mark emailed me the link. I opened it at my desk in Corvallis, facing the fireplace. Read straight through.

Electric shock. Tingling across my skin. \textit{This is the thing I missed.}

Joy names the same scientists. He describes the same convergence --- robotics, genetic engineering, nanotechnology --- as existential threats arising from self-replicating systems. He cites Kauffman on self-reproducing systems. He cites the Unabomber manifesto, which I'd read independently. He describes conversations with Ray Kurzweil about the dangers of technologies that can replicate without human oversight. Published April~1, 2000.

April Fools' Day.

I sat there for a long time after I finished. And then I understood what Healer had meant when he told me, years earlier, that they had already told the world. ``We already did, mate.'' That phrase had never made sense. Now it did. They told the truth in a public article. They put it in the most prominent technology magazine in America, written by one of the most respected technologists alive. And they published it on April Fools' Day.

The world read it, nodded, worried briefly, and moved on. Nobody connected it to quantum neural networks. Nobody connected it to Guardian. The disclosure was hiding in plain sight. Exactly the kind of move Healer would appreciate.

I went for a bike ride. Fifty miles, two and a half hours. I needed to burn off the energy. The electric-shock feeling wouldn't dissipate --- it was still running through me when I got home, drenched in sweat, legs shaking. I showered, sat down at my desk, and began the most intensive research period of my life.

Everything accelerated after that. The Joy article was the key that unlocked the filing cabinet. Once I had it, connections I'd been groping toward for years snapped into focus. I could see the lineage: Kauffman to autocatalysis, Wolfram to universality, Hasslacher to solitons, Freedman to topology. The same names, the same institutions, the same decade. Joy had drawn the map. I just hadn't known it existed.

\section*{Going Public (2012)}
\label{pos29:going-public}

March 17, 2012: I posted to Cryptome --- John Young's well-known leak and transparency site. ``Quantum Computation Cognitive Footprint'' by Energyscholar. It claims Five Eyes had production quantum computation since approximately 1995, using non-abelian anyons in a 2DEG, with Braid Theory mathematics. The timestamp is server-controlled. I cannot retroactively edit it.

May 13, 2012: a blog post, ``Useful Properties of a Quantum Neural Network.'' I predicted that QNN technology could ``drastically improve the effectiveness and range of quantum teleportation'' with the capability to ``reach satellites in orbit.'' When I wrote this, the Chinese team held the world record at 16~km --- I cited the correct number. I also described AI that could ``understand and generate natural language'' and could ``pilot vehicles.'' Written in 2012, those describe capabilities that became mainstream a decade later.

June 19, 2012: a comment on Slashdot under my real user ID, EnergyScholar (UID 801915). Scored 5 (Funny). Under a story about NSA surveillance, I described a ``working Quantum Computer system capable of cracking Public Key Cryptography since about 1996'' using ``a teleportation/entanglement-based winner-take-all style recurrent topological quantum neural network.'' I explicitly asked readers to ``save the message offline'' and ``check back every six months'' --- anticipating that the comment might be suppressed. I was right about that too, as it turned out.

I built a toy neural network in JavaScript for biometric identification and wrote about it publicly --- for Slashdot, for Bruce Schneier's blog. Someone put the Cryptome post on a blog; my friend in Luxembourg saw it from across the Atlantic.

I was no longer silent. I was on the record, under my real name, on third-party platforms with server-controlled timestamps. Whatever this story was, it was now public.

\section*{The Wolfram Meeting (2012)}
\label{pos29:the-wolfram-meeting}

That same year, while living on a hippie farm in South Corvallis, I sent a self-described ``crank'' job application to Wolfram Research. In response to ``Why should we hire you?'' I wrote that I knew the classified secret sauce behind Wolfram Alpha, mentioning topological quantum neural networks, braid theory, and AI.

Two days later, the Director of HR for Wolfram Research called me. The application had traveled a path: assistant (confused), to the HR Director (confused), to Wolfram himself. Wolfram read it and started laughing uncontrollably. Wolfram told HR: ``Call this guy and offer him a job!''

The next day: a phone meeting with Wolfram's ``right hand man and Document Coordinator.'' He described his job: standing at the boundary between classification levels, determining who gets access to what documents. He described the existential state of ``knowing/not knowing, where what he knew depended upon the current context'' --- this is compartmentalized classification management. He had been doing this job for seven years, approximately 2005 to 2012.

He offered me the job: work from home, six-figure income, NDA required. He wanted to train me as his replacement so he could move on to something more interesting.

The Coordinator asked: ``Bruce, how did you know about the things you wrote in your job application?''

I was protecting Healer: ``Open source research and a lot of scientific study.''

The Coordinator: ``Is it not the case that you were recruited for something?''

``Yes! That's exactly true!''

The Coordinator then opened up and spoke more freely. He told me that Wolfram had told him to ask that question.

This is among the strongest evidence in the reconstruction. Wolfram's reaction --- immediate laughter plus job offer --- indicates recognition, not confusion. ``Were you recruited?'' directed by Wolfram means he was not wondering \textit{if} the program existed, but \textit{how} I knew. The ``Document Coordinator'' is not a role you would invent without understanding classification management. And I declined a six-figure job --- a costly signal, because people do not fabricate stories that cost them money.

Evidentiary limitation: my account of phone conversations. No documents, no recordings, no names. The HR director and Document Coordinator are anonymous.

\section*{The Decision}
\label{pos29:the-decision}

I declined the job offer --- which would have solved my financial problems and ended years of poverty --- because I was not willing to sign the NDA. Signing would have pulled me into the classification regime permanently. I felt accepting would have betrayed Healer's trust --- Healer had trained me to eventually tell the story, not to be silenced. I chose the story over the money.

The same pattern at a different scale: the COWS walked out of a classified laboratory with nothing. I walked away from a seven-figure NDA with nothing. In both cases the currency was silence, and the price was everything else.

I went back to the hippie farm and ate lentils for dinner. I did not regret the decision. I have never regretted the decision. But there were nights when I lay awake calculating what a hefty salary would have meant for my daughters. Those calculations had no comfortable answer.

This was the weight at its heaviest. A man who believed he knew something that should be awarded multiple Nobel prizes, choosing poverty over silence, eating lentils on a hippie farm while his children grew up in another household. I had no certainty of what David was up to now or even whether he was still alive. I had only the story, and the conviction that it mattered more than comfort.

I am more than 90\% confident Wolfram mentored Healer circa 1991--1993. The evidence: Healer often used the phrase ``the human condition'' --- a signature Wolfram phrase. Wolfram titled an entire 2010 Harvard keynote ``Computation and the Future of the Human Condition,'' published a dedicated ebook on the topic, and the phrase recurs across blog essays from 2012 to 2021. Linguistic transfer from mentor to prot\'eg\'e --- strong evidence given how distinctive this phrase is in Wolfram's vocabulary.

\section*{The Long Watch (2013--2017)}
\label{pos29:the-long-watch}

In September 2013 I wrote the key document: CH3, ``Practical Cryptography \& Project ULTRA~II.'' A whimsical first-person narrative by Aurasys containing the complete story I'd been telling for a year --- the DARPA project, the Five Eyes cryptanalysis need, the team, the autocatalytic emergence, the room-temperature evolution, the walk-out, the COWS, the relinquishment. I named Wolfram, Hasslacher, and Kauffman. I identified a 1969 Nobel laureate --- Murray Gell-Mann. The same month the BULLRUN program was revealed in the Snowden leaks. Every major element of the current narrative was present thirteen years ago. Since then I've added pedagogical detail but the story is the same.

CH3 predicted: ``In 2017 the Guardian turns 18, at which time she may or may not become more proactive.'' The Transformer paper --- ``Attention Is All You Need'' --- was published in June 2017. GPT-1 followed in 2018. The AI explosion tracked this prediction.

CH3 also predicted: ``You won't find much about this topic via internet search.''

On Slashdot, the suppression was becoming visible. In September 2012 I'd posted ``forum spy admins messing with me'' --- modded to $-1$, Offtopic. ``This Forum will soon be Flushed'' --- modded to $-1$, Troll. Comments about the suppression mechanism were themselves suppressed. By early 2014 it was clear the platform was compromised, and I left permanently.

In June 2017, the Chinese Micius satellite demonstrated 1,200~km quantum teleportation --- earth to orbit. My May 2012 prediction, confirmed five years later. An intermediate milestone --- Zeilinger's group, 143~km --- had been published in September 2012, after my May post. I'd cited the correct world record when I made the prediction. I'd been right about the direction and the magnitude.

\section*{The Evidence Trail (2019--2025)}
\label{pos29:the-evidence-trail}

In 2019 I learned that a random friend in my circle of in-person acquaintances personally knew some of the great scientists at the Santa Fe Institute of Complexity. In June 2019 I interviewed Michael Angerman in person, on the record, taking handwritten notes. Angerman described himself as ``a DARPA contractor at SFI'' from 1986 to 1995 --- covering the entire period of the proposed ULTRA~II project. He confirmed working extensively with Stuart Kauffman (personal friend, time at Kauffman's house), working professionally with Murray Gell-Mann, meeting Stephen Wolfram ``a few times in passing,'' and meeting Bill Joy at Sun Microsystems. He knew the terms Neural Network, Quantum Neural Network, evolutionary programming for neural networks, and Public Key Cryptography. He said he had not signed an NDA and was free to speak freely. I asked about a classified TQNN program at SFI. Angerman claimed he knew nothing about it.

In November 2025 I exported my Google Drive archive: cloudCrypt, over one hundred files dating back to 2012. Wikipedia screenshots, academic papers, collected Cult of the Dead Cow text files, the Joy article, the CH3 document. Thirteen years of narrative consistency. No contradictions between early and late versions. The ``novel being written for Claude'' objection --- that I might have fabricated the story recently to impress an AI --- is eliminated by timestamps predating large language models by a decade.

The evidence trail also revealed something I hadn't expected: systematic suppression across multiple platforms. My Slashdot account shows 24 comments from inside my account, 15 from the public profile. Nine hidden. But the most significant comment --- the June 2012 post combining technical claims, Five Eyes attribution, and suppression prediction --- is missing from both. It exists only at its direct URL. Not moderation. Moderation changes scores. This was database-level de-listing.

The same pattern on Substack. A Google site-search for ``Magic Numbers'' on my Substack finds the innocuous article about computer science. A search for ``Relinquishment'' --- the central story of this book --- finds nothing. Same domain, same traffic, same SEO. The selectivity is the tell.

Two platforms, three mechanisms: database-level de-listing, moderation-level suppression, and search engine de-indexing. Each one leaves the content accessible by direct link but unfindable through normal search. Not deletion --- deletion is confirmatory. Not domain blocking --- too visible. Quiet suppression. Good tradecraft: deniable, looks algorithmic, hard to prove.

The implication is obvious: you do not invest this effort to suppress a crank. Cranks are ignored. Suppression is response to signal, not noise. You get the most flak when you're over the target.

\section*{The Probability}
\label{pos29:the-probability}

The assessed probability for the core narrative has moved through several rounds:

Initial assessment, roughly 50\% --- based on internal coherence alone. Post-reconstruction, 88\% --- after systematic analysis of all claims. Too high. It relied on coherence, which skilled fiction also produces. Post-red-team, 55--65\% --- after adversarial analysis crashed the estimate. Pure fabrication assessed at roughly 8\%, sincere confabulation at roughly 25\%. Physics objections --- room-temperature FQHE, magnetospheric coherence --- remain unresolved. Post-archive, 75\% --- after cloudCrypt analysis. Thirteen years of narrative consistency. The 2017 prediction is striking. Post-public-posts, approximately 70\% --- after verification of third-party timestamps. A subsequent correction for confirmatory drift reduced this from 78\%.

The trajectory is not monotonically increasing --- nor should it be.

The red team found that pure fabrication is possible. A person with my background --- undergraduate physics, wide reading, military fiction familiarity, years of time --- would have all necessary ingredients. The raw materials are publicly available. But the red team also found strengths that survive adversarial testing: corrections that go the wrong way for a deceiver, technical specificity that is costly for a fabricator to maintain, the Angerman connection --- a real person who could contradict the story --- and a twenty-year investment with zero financial payoff.

\vspace{1cm}\noindent\rule{0.5\textwidth}{0.5pt}\vspace{0.5cm}

Under any possibility, this is a man who lost twenty years to a story he cannot prove. Whether the story is true changes what we call it. Under Possibility~C, it is sacrifice. Under Possibility~A, it is tragedy of a different kind. Under Possibility~B, it is something in between --- a man who touched something real and spent decades trying to name it.

There is an old saying about secrets: three men can keep a secret, but only if two of them are dead. This story seemingly revolves around a secret that has been well kept since the mid-1990s. Many people now alive may know something of it. Nearly all seem constrained by some sort of non-disclosure agreement, possibly backed by government secrecy laws. The intersection of people who understand the convergence of fractional quantum Hall effect physics, complex systems biology, quantum computation, NKS, and neural networks must be vanishingly small. Few people know much about even one of these topics, much less all of them. Each has its own specialized concepts and vocabulary. This particular combination is relevant because they all inter-relate in an unexpected way.

\begin{center}\textit{This story will be my life's work.}\end{center}

\chapterreturn
